\paragraph{The finite sum on a class pair.}
\begin{proposition}[Evaluation of the finite kernel on a pair of classes]\label{prop:unserved-reader-integral}
At a level-free place of Lemma~\ref{hyp:unserved-level-free}, the value of the finite kernel on this pair of classes is the sum
\eqref{eq:unserved-reader}. For each coefficient \(a\) and translation point \(P\) occurring there,
\[
 v_w(a)\ge0,\qquad \gcd(\operatorname{ord}P,p)=1,
\]
  and
the Taylor coefficients in the prescribed differential frames
can be transported coefficientwise by the embedding at \(w\).
\end{proposition}

\begin{proof}
\proofstep{Use the local levels of the two self-dual companions.}

Fix both classes of the matrix entry. Write \(L,L'\) for their
companion lattices in the two standard Hermitian spaces. They are
self-dual by the localization argument of
Proposition~\ref{prop:incident-companion-lattices}. Let \(\Lambda\) be the
embedded period lattice obtained from the inverse half different
times \(\operatorname{Hom}_{\mathcal O_E}(L,L')\), with the
curve period scale. At the finite test places the two half lattices
\(A,B\) have opposite level exponents.

If \(c\) is the positive
integer defining the level of the finite kernel, its prime divisors
are in \(S=\{T_u,g_0,g_1\}\cup Q(\mu)\), and
\[
 cA+cB\subseteq\Lambda,\qquad c\Lambda\subseteq A\cap B.
\]
These follow by comparing each of the exponents with the
valuation of \(c\) locally, including the zero exponents off the
level. Write \(S_X,S_A,S_B\) for representatives, respectively, of
\[
 A/(A\cap\Lambda),\quad
 (\Lambda\cap A)/(\Lambda\cap A\cap B),\quad
 B/(B\cap\Lambda).
\]
In particular \(h=|S_A|\) is prime to \(p\). Indeed the middle
quotient is finite and killed by \(c\).

\proofstep{Compute the finite Fourier expansion.}

Put \(E_H(x,y)=\Im H(x,y)\) for the alternating period form,
\(\alpha(x)=\exp(\pi i\sqrt d\, H(x,x)/2)\), and let
\(\Xi\) be the actual finite test extended by zero off the
embedded rational Hom space. For the normalized theta
\(\Theta_{L,L'}\), the finite part of the entry is
\begin{equation}\label{eq:unserved-reader}
 \sum_{x\in S_X,\ y\in S_B}
   \alpha(x)\alpha(y)e^{-\pi i E_H(x,y)}
   \left(\frac1h\sum_{s\in S_A}
      \Xi(s-x)e^{2\pi i E_H(s,y)}\right) U_{x+y}\Theta_{L,L'} .
\end{equation}
The Fourier sum at \(y\) uses the inverse commutator \(\exp(-2\pi i E_H(s,y))\), with the counting
measure divided by \(h\). The kernel reading uses volume one.
Applying \(U_x\) to the resulting \(U_y\) gives the displayed
Weyl factor. Since \(E_H(\Lambda,\Lambda)\subset\mathbb Z\)
and \(\sqrt d\,H(v,v)/2\in\mathbb Z\) on \(\Lambda\), all three
phase factors and the Fourier character are roots of unity.
The translates have order dividing \(c\) modulo \(\Lambda\).

\proofstep{Check the additive and determinant character factors.}

The word support is
a characteristic function. The type tests and twist tests on
a residue matrix multiply finite additive characters and values
of finite multiplicative characters, with zero on the excluded
residue classes. Each nonzero value is a root of unity.
Thus, for a nonzero finite character factor \(\xi\) of order \(n\),
\[
  n\,v_w(\xi)=v_w(\xi^n)=v_w(1)=0 .
\]
 In a
moved test the extra factors are integer powers of square roots
of the residue cardinality and local determinant characters.
At a type place the two movers preserve the submitted lattices,
so their determinants are local units and the residue characters
suffice.

At a twist place the uniformizer value and its inverse
are integral away from that place's rational prime and the type
levels. Indeed raise the ideal of the place to a principal power
with generator congruent to one at the character conductor.
The defining algebraic infinity type expresses that power of
the value as a product of conjugates of the generator to
integer powers, up to a root of unity. If \(\mathfrak P^A=(\alpha)\) with \(\alpha\equiv1\)
at the conductor, the value \(u\) satisfies

\[
  u^A=\zeta\prod_\tau\tau(\alpha)^{a_\tau}
\]
in the algebraic normalization. Thus
\[
       A v_w(u)=\sum_\tau a_\tau v_w(\tau(\alpha))=0
\]
for \(w\) outside that rational prime and the type levels. The off-type matching of the
column character with this ideal character gives the assertion
for the actual local value. The types are disjoint from the
auxiliary twist places. Finally the pinned terminal row either
uses an indicator or its finite Fourier reading. Its nonzero
values are integral additive character values, with a possible
factor \(q^{-Mn}\) at a place of cardinality \(q\), rank \(n\)
and depth \(M\). At the coefficient place under consideration, \(q\) is a unit. For any nonzero additive-character value \(\xi\) this gives
\[
  v_w(q^{-Mn}\xi)=-Mn\,v_w(q)+v_w(\xi)=0.
\]
 The
nonterminal classes use the matrix test, by the terminal pinning
of the supplementary rows. This proves integrality of every
scalar multiplying a translate in \eqref{eq:unserved-reader}.
It applies also when an entire entry vanishes.
\label{cor:unserved-kernel-charts-1}

All Taylor identities here are identities of coefficients.
\proofstep{Transport the Taylor coefficients in their fixed normalization.}

For clarity, algebraicity of the translated derivatives can
also be tested at one auxiliary good prime. Choose it away from
the denominators of the period comparison, the translate and a
rational basis of the Hom directions, away from the discriminant,
and split completely in the host. The fourfold expansion and the
baseline derivative estimate of Section~\ref{subsec:baseline-jets}
apply for every isomorphism of \(\mathbb C\) with \(\mathbb C_p\)
over a fixed embedding of \(\overline{\mathbb Q}\) at this auxiliary
prime. For each fixed derivative \(c\) in those directions the resulting bound \(M_c\) is uniform in the chosen isomorphism:
\[
  |\iota(c)|_p\le M_c\qquad\bigl(\iota:\mathbb C\xrightarrow{\sim}\mathbb C_p\text{ over the fixed algebraic embedding}\bigr).
\]
 In particular, the argument matrices are
integral after a fixed denominator clear for the rational
directions. The embedded matrices of a rational basis over the
quadratic field span the complexified Hom at the CM type, by
separability of \(E/K\).

For a fixed total order \(d\), the induced action on derivatives
is the degree-\(d\) tensor of the inverse direction matrix. Its
entries are algebraic and fixed. Thus every labelled coefficient
is a fixed finite algebraic linear combination of the bounded
directional derivatives. The triangle inequality gives the same
uniformity in the field isomorphisms.
Suppose one coefficient \(c\) were transcendental, and fix one
of the comparison isomorphisms \(\iota\colon\mathbb C\to\mathbb C_p\).
For each \(n\ge1\), the map 
\[
  c\mapsto p^{-n}c
\]
 extends by
transcendence bases to an automorphism \(\alpha_n\) over
\(\overline{\mathbb Q}\). The same directional bound would then
apply to all the isomorphisms \(\iota\circ\alpha_n\), whereas
\[
 |(\iota\circ\alpha_n)(c)|_p=p^n|\iota(c)|_p\longrightarrow\infty .
\] Division by the multi-index factorial preserves
algebraicity. The finite sums, linear substitutions and the
map on derivatives can consequently be transported by the
specified coefficient embedding at the original \(w\).

\end{proof}

\paragraph{Removing the differential generator in the actual frames.}

A differential generator can change the valuation in a scalar chart. In the actual integral frames the following comparison removes that factor.

\begin{proposition}[Removing the generator in integral frames]\label{prop:unserved-normalized-frame}
At the enhanced places in this comparison, including \(p=d\),
one can choose the class frames so that the substitution of \eqref{eq:unserved-normalized-frame} reads
\[
 Z_\tau\longmapsto \rho_\tau(P_{i+1})Z_\tau\rho_\tau(P_i)^{-1}.
\]
 The axis
generators and the resulting scalar factors in the coefficient
tests are units in the chosen coefficient valuation.
\end{proposition}

\begin{proof}
\proofstep{Choose simultaneous rational unitary frames.}

For each stage \(i\) and each of its classes choose a rational
unitary matrix \(P_i\) whose columns span its lattice at all
places above \(p\). Such a matrix exists simultaneously there.
The selected host extension \(E/F\) is unramified at finite places.
At a split place the two summands of a self-dual lattice are
mutually dual.
In paired local coordinates the Gram block between dual bases \(e_i,e'_j\) is
\[
  \big(h(e_i,e'_j)\big)_{ij}=I .
\]

At an unramified nonsplit place, orthogonal
splitting and surjectivity of the unramified unit norm identify
two unimodular Hermitian lattices. The resulting local
transporters are open. Multiplication by nearby scalar norm-one
elements avoids the eigenvalue \(-1\) at each place. Approximate
their inverse Cayley transforms in the rational vector space of
skew-Hermitian matrices by weak approximation over \(F\), and
apply the Cayley transform back. Thus both the columns and the
inverse coordinates clear into the required lattices by
integers prime to \(p\). The same matrix \(P\) works at every label over the place. Its inverse adjoint property is
\[
  P=(P^*)^{-1},\qquad P^*P=I.
\]

\proofstep{Separate the residual axes.}

For a CM type embedding \(\rho_\tau\) at \(w\) write \(d_\tau\)
for the image of the generator of its axis in the module
spanned by \((\rho_z a)_z\), \(a\) in the inverse half different.
At these enhanced places that ideal is the unit ideal locally.
The entries are integral. For two distinct type labels \(\tau,\tau'\), one can consequently separate their residue embeddings by an integer:
\[
  a\in{\cal O}_E,\qquad \overline{\rho_\tau(a)}\ne\overline{\rho_{\tau'}(a)}.
\]
 Indeed the labels
all extend the same embedding of \(K\), and the relative inertia
in \(E/K\) is trivial at these primes.

Distinct primes are
separated by the Chinese remainder theorem and distinct
embeddings of the same unramified local field are separated
on the residue field. This also applies at \(p=d\), where
\(K/\mathbb Q\) itself is ramified. Acting on the embedded
all-ones vector with the product, over \(z\ne\tau\), of
\(a_z-\rho_z(a_z)\), where
\(\rho_\tau(a_z)-\rho_z(a_z)\) is a unit, isolates the
\(\tau\)-axis with unit coordinate. The module is stable under
these actions. Its axis generator is therefore a unit, and
membership of the generator supplies a finite expression
\[
 d_\tau {\bf1}_\tau=\sum_l r_l(\rho_z a_l)_z,
 \qquad |r_l|_w\le1,\quad a_l\in\mathfrak b^{-1},
\]
where \(\mathfrak b\) denotes the half different.
This last expression keeps its global period directions.

\proofstep{Absorb the unit without changing the norm tests.}

Set \(a_\tau=d_\tau^{-1}\) and use the frame
\(a_\tau^i\rho_\tau(P_i)\) at stage \(i\). In the incident
Hom variable the factor
\(d_\tau a_\tau^{i+1}/a_\tau^i\) is exactly one. The substitution
on the series in \eqref{eq:unserved-reader} is now
\begin{equation}\label{eq:unserved-normalized-frame}
 Z_\tau\longmapsto
    \rho_\tau(P_{i+1}) Z_\tau\rho_\tau(P_i)^{-1},
\end{equation}
with the row of \(Z_\tau\) indexed by the target. These are
admissible choices for the true columns.

The product of the
column minors of a diagram \(\nu_\tau\) gains
\((a_\tau^i)^{|\nu_\tau|}\) on input frames. A determinant twist
of exponent \(s\) gains \((a_\tau^{iN_i})^s\), including signed
exponents. The target formulas use \(i+1\). Every such factor
is a unit, and homogeneity gives the same power on every term
of the corresponding polynomial reading. Thus the original
coefficient norm tests survive this change of frames, while the
Taylor map itself has the precise coefficient in
\eqref{eq:unserved-normalized-frame}. \label{cor:unserved-kernel-charts-2}

\end{proof}

\paragraph{One curve model before the inert orbit estimates.}
\begin{proposition}[One inert curve model for all translates]\label{prop:unserved-common-inert-model}
At the selected inert places \(p\ge5\) with \(p\ne d\), the two classes and all their prime-to-\(p\) translates use one good curve law, tangent unit \(c_w\) and quadratic correction \(q\). In good curve coordinates with logarithm \(l\), the parameter change for all these charts is
\[
 t(X)=l^{-1}(c_w L(X)),\qquad c_w\in{\cal O}_w^\times.
\]

In the even and signal cases the resulting normal form is
\textup{\eqref{eq:tag-i4}}. The odd case uses the doubled orbit and
the zero-step case uses the repeated-variable one-letter
phase. The phase and its coefficient estimates are uniform
in the labels, although the integral series may depend
on the classes and on every variable.
\end{proposition}

\begin{proof}
\proofstep{Fix the model and the common unit.}

The four cases treated here share the following hypotheses on the residue characteristic
\[
  p\ge5,\qquad p\ne d,\qquad \left(\frac{-d}{p}\right)=-1.
\]
Apply Proposition~\ref{prop:good-inert-cm-model} to the discriminant-normalized
period lattice. Changing the integral ideal representative of the
curve prime to \(p\) changes the discriminant root by the inverse
homothety and leaves its scaled lattice fixed. The good-model
homothety is a unit since its twelfth power compares two unit
discriminants. Write \(l\) for the good curve logarithm and \(L\) for the comparison logarithm. The integral parameter change \(t\) fixes the tangent unit through
\[
  l(t(X))=c_wL(X),\qquad t'(0)=c_w\in{\cal O}_w^\times.
\]
Use this one choice of parameter, Frobenius and curve phase \(q\) throughout the place. These precede the choice of the two classes, their ranks,
their diagrams, and the translations in the reader.

\proofstep{Choose the underlying Hom frame before the translates.}

Here is a rational orthonormal chart for the Hom of the actual
pair. For the standard lattices one matrix coordinate carries
the trace Hermitian form on the inverse half different over \(K\).
It is unimodular at \(p\) by duality of the trace lattice
(the half different times its conjugate is the relative
different).
Split off a vector of unit norm from its reduction.
For a lift \(u\) with unit norm, let \(\ell\) be its Hermitian pairing in the linear variable. The integral projection and its kernel are
\[
  x\longmapsto \ell(x)\ell(u)^{-1}u,\qquad \ell(u)\in{\cal O}^{\times},\qquad \ker\ell=u^\perp.
\]
 Such a
vector exists in every nonzero nonsingular Hermitian space over
the quadratic residue field, by polarization. Lift it to the
localized rational lattice.

Orthogonal projection onto its
complement divides only by that unit, so iteration gives an
orthogonal basis with positive rational norms \(n_j/e_j\) whose
numerators and denominators are prime to \(p\). Positivity and
rationality follow from definiteness and the fixed field of
conjugation on \(K\). In four copies of each line choose four
integer squares summing to \(n_je_j\). If \(Q_j\) is their quaternion matrix, then
\[
 Q_j^tQ_j=n_je_j I_4,\qquad
 (Q_j/n_j)^t\bigl((n_j/e_j)I_4\bigr)(Q_j/n_j)=I_4 .
\]
Both \(Q_j/n_j\) and its inverse are integral at \(p\). Use these four copies in each
\(2\)-by-\(2\) block of matrix entries, possible since both
ranks are even. Transport the resulting basis by

\[
  X\mapsto P_{i+1} X P_i^{-1}.
\]
 This is a rational isometry
and spans the Hom lattice over the local quadratic integers.
The basis is chosen before any translate.

\proofstep{Descend the corrected integral section from the fourth tensor.}

On four copies the normalized theta expansion of
Section~\ref{subsec:theta-normalization} has integral coefficients
in the locally matched frame. Its arguments are integral
unitary transforms of the orthonormal directions just
constructed. The one-variable translated section multiplied by
\(\exp(q)\) is integral in group coordinates. For the sum of
phases a unitary change preserves its class modulo logarithms
of principal integral units by the polarized CM norm relation
\eqref{eq:cm-norm-relation}.

Apply the change and extract four
equal coefficient indices on disjoint copies. This proves, in the
orthonormal basis \(b_j\), that the product
\[
  \exp\!\left(\sum_j q(X_j)\right)
  \Theta\!\left(\left(c_w\sum_j\rho_\tau(b_j)L(X_j)
                        \right)_\tau\right)
\]
is integral, for each of the normalized translates \(\Theta\)
in the finite reader. It holds also for an integral linear
combination since the multiplier and arguments are common. A pair
of lattices agreeing only locally at \(p\) gives the same conclusion.
Two integers prime to \(p\) clear the source and target lattice
inclusions in both directions, preserving these local arguments
and prime-to-\(p\) torsion.

\proofstep{Read the phases in the trace-dual cycle variables.}

In an even orbit expand the two parities in a basis and its
trace dual over the quadratic local field. For one pair of parity vectors \(e,o\) in a matrix coordinate, the Gram block is
\[
  \begin{pmatrix}h(e,e)&h(e,o)\\ h(o,e)&h(o,o)\end{pmatrix}=\begin{pmatrix}0&1\\1&0\end{pmatrix}.
\]
 This follows by multiplying the unitary
source and target changes in
\eqref{eq:unserved-normalized-frame} by their inverse adjoints
on the conjugate coordinates. Thus the change from the
orthonormal variables gives the hyperbolic substitution of
Section~\ref{subsec:curve-exponents}, with the trace-dual sums
there called \(E\) and \(O\). The root-of-unity average over
\(p^{d_p}-1\) elements is integral and replaces the summed
polarization by \(B_{d_p}\) modulo logarithms of integral
units.
The averaging denominator is a unit, since
\[
  p^{d_p}-1\equiv-1\pmod p,\qquad v_p\big((p^{d_p}-1)^{-1}\big)=0.
\]

Returning to additive coordinates gives \eqref{eq:tag-i4} of
Section~\ref{subsec:incident-sections}, with exponent
\(-B_{d_p}\) on each entry and argument scale given by the
common tangent unit. This works for any choice of the orbit
representatives and includes the signal orbits. The same
orthonormal identity at odd order is used below with the
trace-dual variables on the doubled orbit. At order one the
labelled basis is already unitary over the quadratic integers.
The norm relation transports the one-variable gauge, and its
polarization at the repeated variable gives the one-letter
phase with the factor \(1/2\). Thus in each case the integral
factor can depend on the classes and on every variable, while
the phase and its coefficient estimates are uniform in both
classes and all entries of their diagrams. \label{cor:unserved-kernel-charts-3}
\end{proof}
